\section{Nostalgia and Death}

\begin{quotex}
What a shame that one must give up life. I was finally getting used to it.
\flright{\textsc{Charles Peguy}}

If you do not meditate on death upon waking in the morning, your entire morning will be wasted; if you do not meditate on death at noon, your entire afternoon will be wasted; and if you do not meditate on death in the evening, your entire night will be wasted.
\flright{\textit{Tibetan Buddhist teaching}}
\end{quotex}


\paragraph{Thanks for the Memories}


\begin{quotex}
The truth is that memory does not consist in a regression from the present to the past, but on the contrary in a progress from the past to the present.
\flright{\textsc{Henri Bergson}, \emph{Matter and Memory}}
\end{quotex}


I am letting go of memories. Why do the oddest events stand out? Seemingly insignificant past events create vivid memories. What was so important 20 or 30 years ago is meaningless today. Do I want that an ancient event risks salvation?

The future is a horizon that we cannot see beyond. We approach it. Younger men look forward to it. I see death, an ugly death. Lying in a hospital bed, hooked up to tubes, a ventilator, and electronic equipment. First the kidneys fail, then gradually other organs, one by one, until the end.

Nowadays, death is expected to be safe, legal, and rare. It is not morbid to envision one's death, but rather it enhances life and offers perspective.

\paragraph{Maya}


\begin{quotex}
If the soul's mystical union with God is forgetting of the phenomenal world and recall of God, death is simultaneously the call from above and forgetting below.
\flright{\textsc{Valentin Tomberg}, \emph{Meditations on the Tarot}}
\end{quotex}


Maya is the world of phenomena. It is called ``illusion'', not because our senses are mistaken, but because we assume it is the real world. That is because Maya is persistent and forcibly intrudes herself into our consciousness. Sights, sounds, smells enter our minds uninvited. Are they good, evil, neutral? Do they evoke fear, hunger, desire, lust? They force a decision: Attack, retreat, hide, ignore, or perhaps enjoy.

Despite efforts at detachment, I find I am still deeply in love with Maya. The world of sensory experience can also be beautiful, delightful, enticing. Children are thrilled with each new experience, and want to repeat it. But at what point does repetition no longer produce the desired effect? Or, more likely, we will eventually learn to see that our judgments of this world were all too often far off the mark.

If you want to detach from Maya, then try Freddie's suggestion:

\emph{You go to women? Do not forget the whip!}


\paragraph{The Call from Above}


Since the world of Maya is so dominant, it is impossible to ignore. In contrast, the Call from God may be faint and is often missed. The first task is to silence the mind, which is difficult given our attachment to Maya. Even if there has been some success in that task, we are then faced with a more subtle problem. The inner life is still in a turmoil, not from Maya, but is self-caused. There are negative emotions, useless fantasies, unfulfilled desires, and so on. These need to be overcome, in turn. Tomberg explains what happens next:

\begin{quotex}
We must renounce the joy of living in the richly filled images that arise before one's inner vision and obliterate the imaginations that have been reached. In their place darkness and emptiness must be created. And in this darkness and emptiness one must immerse oneself with one's entire soul.
\end{quotex}


That does not sound so great; moreover, it does not feel so great. Tomberg goes on:

\begin{quotex}
The soul feels itself in complete and absolute inner loneliness. One experiences oneself as if forsaken by all human beings and by the gods.
\end{quotex}


This emptying process needs to continue, until the soul learns to hear in the spiritual world. This process actually describes the traditional states of purification, illumination, and union.

\paragraph{Inner life is opaque}


\textbf{St. Augustine} asserts that the truth is found in our interiority. Unfortunately, he also shows that our inner life is opaque to us; and \emph{a fortiori}, the inner lives of others are likewise opaque. Nevertheless, it is all too common in public discourse to claim to know and understand the hidden motivations of others, an easy way to dismiss them and their ideas.

In our particular judgment, our hidden motivations will be revealed to us; only then, will we fully understand God's justice. In the general judgment of all, the inner lives of all people will be revealed. Once, God's justice will be full known and accepted. In our current situation, it is difficult to see God's justice, which explains the depth of the resistance to this understanding.

\paragraph{Resurrection Body}


There are four qualities of the resurrection body:

\begin{enumerate}


\item \textbf{Impassibility}: It does not suffer pain or death

\item \textbf{Clarity}: It is free from deformity and is filled with beauty and radiance

\item \textbf{Agility}: The body obeys the soul and can instantly travel from place to place

\item \textbf{Subtlety}: It has a spiritual nature and mastery over the physical. E.g., it can pass through walls.
\end{enumerate}


One of the most important insights of \textbf{Wolfgang Smith} is the distinction between the corporeal plane and the physical plane. Hence, the resurrection body is \emph{corporeal}, i.e., it has the qualities and appearance of a body, but no longer \emph{physical}, which is associated with quantity.

\paragraph{Gethsemane}


Jesus finds himself alone at Gethsemane, lost in God's silence on the one hand, and the sleeping disciples on the other. His friends, one by one, prefer sleep to consciousness. And in the most important spiritual decisions, God is silent. That is because a man must choice freely to follow God's will, without coercion.

\paragraph{Washing the Feet}


In \textit{Inner Development}, Tomberg describes the effect of meditation:

\begin{quotex}
The general effect of meditation consists in the fact that what is spiritual in a human being descends. That is, the superconsciousness -- the higher I -- descends in the human personality, and the angel sends down the stream of illumination in the form of ``washing the feet''.
\end{quotex}


This is understood as the spiritualization of the lowest parts of the human. The incarnation, the descent into hell, demonstrate that. This can be contrasted with older systems, such as Vedanta, Neoplatonism, or Gnosticism which sought instead to escape the material world and merge into the Absolute.

This descent of the superconscious is the opposite of the ascent of Kundalini.

\textbf{Rene Guenon} describes deliverance this way:


\begin{quotex}
If the passage to certain superior states in some way constitutes a progress toward `Deliverance' relative to the state taken as a point of departure, it must nevertheless be understood that when the latter is realized, it will always imply a discontinuity with respect to the immediately preceding state of the being that achieves it, and that the discontinuity will be neither more nor less profound whatever this state may be, since in all cases there is between the `undelivered' and the `delivered' being no relation such as exists between the different conditioned states.
\end{quotex}


Yet there is a higher state than annihilation. While not denying the experiences of the yogis, Tomberg describes an even deeper understanding of union with God:

\begin{quotex}
Union with the Divine is not the absorption of being by Divine Being, but rather the experience of the breath of Divine Love, the illumination by Divine Love, and the warmth of Divine Love. The soul which receives this undergoes such a miraculous experience that it cries. In this mystical experience fire meets with FIRE, Then nothing is extinguished in the human personality but, on the contrary, everything is set ablaze. This is the experience of ``legitimate twofoldness'' or \textbf{the union of two separate substances in one sole essence}. The substances remain separate as long as they are bereft of that which is the most precious in all existence: free alliance in love.
\end{quotex}


This union exemplifies God as Being and God as Love. The conditioned states are not annihilated but retain their substance, while at the same time, the essence of the mystic encounters God's essence. Their wills unite:

\begin{quotex}
The work of the Redemption, being that of love, requires the perfect union in love of two wills, distinct and free ---divine will and human will.
\end{quotex}


The person does not become extinct; rather he becomes fully and truly what he is, as conceived eternally in the Mind of God.


\flrightit{Posted on 2020-04-12 by Cologero}

\begin{center}* * *\end{center}

\begin{footnotesize}
\begin{sffamily}
\texttt{Balder on 2022-04-12 at 17:39 said:}

This came as a good synchronicity, since today I heard about the death of my ex-wife's father, who was for long very close to me. I remember thinking again, that the last thing I how I want to die is in a hospital, in shitty diapers, into some terminal illness. Maybe it's that Mars energy I felt today so strongly, but I started to envision more preferably a violent death in battle than the fate of a sickly old man whose quality of life has deteriorated into terminal state God knows for how long; the worst part is that it is usually relatives who wish to keep the elderly as medicated zombies for years. What is this other than an ill response to the fear of death caused by the society who so dearly venerates so called ``life'' in all aspects and wishes to be done away with the reality ans sanctity of death by all possible means.

Pekka Ervast also wrote in Finland that the task of the initiate is to ``build'' an immortal corporeal body; the part where you described Smith's categorization of the physical and corporeal world struck me, and I think now finally understand what Ervast was speaking about. I recently read again some of Smith's work, and I quite like his insights of science in light of traditional metaphysics.

For long time the idea of a ``immortal corporeal body'' sounded ridiculous, since I didn't quite understand this difference between the physical and corporeal universe. Maybe we can understand this also by the reference to the so called etheric or energetic body, the ``real corporeal body''. In the Finnish epic Kalevala there is the so called ``kirjokannel'' which I think is a reference to the five (or seven, a kannel is usually five-stringed) principles and elements making up the etheric body of man and also the so called chakras. Anyway, I think what the old Finns meant by the Kirjokannel is the same thing as the resurrection body of the Christian tradition.

\end{sffamily}
\end{footnotesize}

